% !PRINTED 668
%% [running head, not body text -- commented out 2026-08-07] Page furniture (running
%% head, folio, signature) is excluded from this edition by standing policy.
%% 668 \hfill \gesperrt{Zweiundzwanzigster Abschnitt.} \hfill \S.\,184.
\weberpage{668}
%<<CONT
worin $\varrho_n$ eine Einheitswurzel vom Grade $m^2$, und $\alpha_n$
eine Zahl in $\Omega_m$ ist, und es w\"are noch zu beweisen, dass
$e(r^n)$ das \gesperrt{Quadrat einer Einheit} in $\Omega_m$ ist.

Da $e(r)$ eine reelle Einheit ist, so ist

\weq{2}{e(r^n) = e(r^{-n}),}

und wir wollen noch festsetzen, dass das Vorzeichen der Wurzel so
bestimmt sei (was wir bei passender Annahme \"uber $\varrho_n$ immer
annehmen k\"onnen), dass

\weq{3}{\sqrt{e(r^n)} = \sqrt{e(r^{-n})}.}

Das Product $(r^n, x_0)\,(r^{-n}, x_0)$ ist nach der Voraussetzung
eine Zahl des K\"orpers $\Omega_m$, die sich durch die Substitution
$(r, r^{-1})$ nicht \"andert, d.\,h. eine reelle Zahl. Ferner ist
nach \S.\,174, (5)

\[
(r^{m_2 n}, \eta)\,(r^{-m_2 n}, \eta) = \pm p,
\]

also gleichfalls reell. Ebenso ist $\alpha_n \alpha_{-n}$ reell, und
daraus ergiebt sich nach (1), dass

\[
\varrho_n \varrho_{-n} = \frac{(r^n, x_0)\,(r^{-n}, x_0)}
{e(r^n)\ \alpha_n \alpha_{-n}\ \Pi(r^{m_2 n}, \eta)\,(r^{-m_2 n}, \eta)}
\]

eine reelle Zahl ist, die, weil sie eine Einheitswurzel ist, nur $=
\pm 1$ sein kann.

Nun k\"onnen wir in der hieraus f\"ur $n = 1$ sich ergebenden
Gleichung

\[
(r, x_0)\,(r^{-1}, x_0) = \pm\ e(r)\ \alpha_1 \alpha_{-1}\ \Pi(\pm p)
\]

jede Substitution $(r, r^n)$ ausf\"uhren, woraus hervorgeht, dass das
Zeichen von $\varrho_n \varrho_{-n}$ von $n$ unabh\"angig ist, dass
also

\weq{4}{\varrho_n \varrho_{-n} = \varrho_1 \varrho_{-1} = \pm 1}

ist. Nun leiten wir aus (1) weiter her:

\weq{5}{\varrho_n \varrho_1^{-1} \sqrt{e(r^n)}\ \sqrt{e(r)}^{-n}
= \frac{(r^n, x_0)\,(r, x_0)^{-n}}
{\alpha_n \alpha_1^{-1}\ \Pi(r^{m_2 n}, \eta)\,(r^{m_2}, \eta)^{-n}},}
% [source sic: (5) prints both \varrho_1^{-1} and \alpha_1^{-1}; formation from (1) and the displayed -n powers requires \varrho_1^{-n} and \alpha_1^{-n}]

und die rechte Seite ist eine Zahl des K\"orpers $\Omega_m$. Die
linke Seite zeigt aber, dass es eine Einheit ist, und wir k\"onnen
demnach, wenn $\mathfrak{E}(r)$ eine reelle Einheit bedeutet,

\weq{6}{\varrho_n \varrho_1^{-1} \sqrt{e(r^n)}\ \sqrt{e(r)}^{-n} = r^k \mathfrak{E}(r)}
% [source sic: (6) repeats the inherited \varrho_1^{-1}; the derivation requires exponent -n]

setzen. Substituiren wir diesen Werth in die Formel (5) und machen
die Substitution $(r, r^{-1})$, so ergiebt sich

\weq{7}{\varrho_{-n} \varrho_{-1}^{-1} \sqrt{e(r^n)}\ \sqrt{e(r)}^{-n} = r^{-k} \mathfrak{E}(r),}
% [source sic: (7) prints \varrho_{-1}^{-1}; the substituted form of the preceding relation requires exponent -n]
%CONT>>
